% =============================================================================
%  Lab 10 - Instruction Decoder I
% =============================================================================

\chapter{Lab 10 --- Instruction Decoder I}

\section{Objective}
% Two or three sentences, your own words.

\section{Specification}

\subsection*{Your four rows of the control table}
% Only the DIFFERENCES from CTRL_DEFAULT. Write "--" where the default stands.

\begin{table}[htbp]
\centering
\caption{Control table, part 1. Differences from \texttt{CTRL\_DEFAULT} only.}
\label{tab:lab10-ctrl}
\small
\begin{tabular}{llllllll}
\toprule
instr & opcode/func & \texttt{ALU\_func} & \texttt{ALU\_Asel} & \texttt{ALU\_Ssel} & \texttt{IR\_Xsel} & \texttt{RF\_strb} & \texttt{RF\_Wsel} \\
\midrule
\texttt{ADD}  & \texttt{000000}/\texttt{100000} & & & & & & \\
\texttt{SLL}  & \texttt{000000}/\texttt{000000} & & & & & & \\
\texttt{ADDI} & \texttt{001000} & & & & & & \\
\texttt{LUI}  & \texttt{001111} & & & & & & \\
\bottomrule
\end{tabular}
\end{table}

\subsection*{The safe default}
% State the values CTRL_DEFAULT gives to RF_strb, DM_en, DM_strb, RF_oe and
% NIAG_sel, and say in one sentence why those particular values are the SAFE
% ones for an instruction the decoder does not recognise.

\section{Design}
% Answer all three:
%   (a) Why must CTRL_DEFAULT be assigned BEFORE the case statement? What
%       exactly does the synthesiser build if you do not, and what would the
%       symptom be in Lab 12?
%   (b) Why must the SPECIAL opcode be decoded on BOTH the opcode and the func
%       field? What happens if you check only the opcode?
%   (c) ADD and SLL leave IR_Xsel and IR_Xlsft at their defaults. Is that a
%       mistake? Explain what a DON'T CARE is and why leaving them alone is
%       correct here.

\section{Implementation}

\vhdlfile{code/lab10/Decoder.vhd}{\texttt{Decoder} --- instruction decoder, part 1}

\section{Verification}
% Vectors, then the coverage argument.
%
% A decoder's coverage argument is different again: the input space is the
% instruction encoding, and the interesting subset is
%     one representative per implemented instruction
%   + one unimplemented opcode
%   + one unimplemented func within SPECIAL
%   + at least one SEQUENCE, to detect latches.
% State that, and say why the sequence test is the one a single-instruction
% testbench cannot make.

\begin{table}[htbp]
\centering
\caption{Test vectors.}
\label{tab:lab10-vectors}
\begin{tabular}{cllp{5.4cm}}
\toprule
\# & instruction & encoding & what it proves \\
\midrule
1 & & & \\
2 & & & \\
3 & & & \\
\bottomrule
\end{tabular}
\end{table}

\section{Results}
% Evidence: E1, E2, E4, E7.

\figplaceholder{the RTL Viewer for \texttt{Decoder} --- combinational logic and
NO storage elements}

% E1: state the error and warning counts, and say EXPLICITLY that there is no
%     inferred-latch warning.
% E4: the register count. It MUST be zero. Say why any non-zero count would mean
%     the design is wrong rather than merely expensive - and name the specific
%     bug it would indicate.
% E7: quote the transcript line showing every check passing, including the
%     order-dependence test.

\section{Discussion}
% REQUIRED. One bug, how you found it, what it cost.
%
% Then answer this explicitly:
%     You have now met "no assignment means hold the previous value" three
%     times: Lab 1, Lab 3 and Lab 10. In which of those did you WANT that
%     behaviour, and what makes this one the most dangerous of the three?
